% Source-aligned English for printed p.158.
\SourcePage{158}

\LayerTarget{tohoku:construction:inverse-image-comparison:p2}
\noindent extends uniquely to a homomorphism of cohomological functors
\[
R\Gamma^Y_\Psi\longrightarrow (R\Gamma^X_\Phi)g.
\]
We have therefore proved:

\LayerTarget{tohoku:prop:3.2.2}
\textsc{Proposition 3.2.2.} --- \textit{Let $f:X\longrightarrow Y$ be a
continuous map, and let $\Phi$ and $\Psi$ be antifilters of closed subsets of
$X$ and $Y$, respectively, such that $B\in\Psi$ implies
$f^{-1}(B)\in\Phi$. For every sheaf of abelian groups $F$ on $Y$, there is a
unique family of homomorphisms}
\[
H^p_\Psi(Y,F)\longrightarrow H^p_\Phi(X,f^{-1}(F))
\qquad(-\infty<p<+\infty)
\]
\textit{forming a homomorphism of cohomological functors whose degree-zero
part is the natural homomorphism.}

These homomorphisms will be called the natural homomorphisms. Their uniqueness
implies an evident transitivity property, whose statement is left to the
reader.

\LayerTarget{tohoku:construction:restriction-homomorphisms}
In particular, if $Y$ is a subset of $X$ and we put
$H^p_\Psi(Y,F)=H^p_\Psi(Y,F\vert Y)$, where $F\vert Y$ denotes the
restriction of the sheaf $F$ to $Y$, there are \textit{restriction
homomorphisms}
\[
H^p_\Phi(X,F)\longrightarrow H^p_{\Phi\cap Y}(Y,F),
\]
where $\Phi\cap Y$ denotes the trace of $\Phi$ on $Y$.

\LayerTarget{tohoku:section:3.3}
\textbf{3.3. Acyclicity criteria.} The developments in this section, which
will be very useful below, are due to Godement and will be treated in detail
in Godement's book [9] cited in the Introduction.

\LayerTarget{tohoku:lemma:3.3.1}
\textsc{Lemma 3.3.1.} --- \textit{Let $F:C\longrightarrow C'$ be a covariant
functor between abelian categories, and suppose that every object of $C$ is
isomorphic to a subobject of an injective object. Let $\mathbf M$ be a class
of objects of $C$ satisfying the following conditions: (i) every $A\in C$
admits a monomorphism into some $M\in\mathbf M$; (ii) every $A\in C$
isomorphic to a direct summand of an $M\in\mathbf M$ belongs to $\mathbf M$;
(iii) for every exact sequence
$0\longrightarrow M'\longrightarrow M\longrightarrow M''\longrightarrow0$
in $C$ with $M',M\in\mathbf M$, one also has $M''\in\mathbf M$, and the
sequence
$0\longrightarrow F(M')\longrightarrow F(M)\longrightarrow
F(M'')\longrightarrow0$ is exact. Then every injective object of $C$ belongs
to $\mathbf M$, and $R^pF(M)=0$ for every $M\in\mathbf M$ and $p>0$.}

First let $I$ be injective in $C$. By (i), $I$ embeds into some
$M\in\mathbf M$ and, since $I$ is injective, it is isomorphic to a direct
summand of $M$. Thus $I\in\mathbf M$ by (ii). Now let $M\in\mathbf M$ and
prove that $R^pF(M)=0$ for $p>0$. Take a right resolution of $M$ by injective
objects $C^i$, $i\geq0$, and let $Z^i$ denote its subobject of cycles in
$C^i$. It is enough to prove that the sequences
\[
\begin{gathered}
0\longrightarrow F(M)\longrightarrow F(C^0)\longrightarrow F(Z^1)
\longrightarrow0,\\
0\longrightarrow F(Z^1)\longrightarrow F(C^1)\longrightarrow F(Z^2)
\longrightarrow0,\quad\text{and so on}
\end{gathered}
\]
are exact. By (iii), it is enough to prove that all $Z^i$ and $C^i$ belong to
$\mathbf M$. This is already known for the $C^i$, since they are injective,
and follows for the $Z^i$ by induction from (iii).

\LayerTarget{tohoku:cor:3.3.1}
\textsc{Corollary.} --- \textit{For every $A\in C$, the groups $R^pF(A)$ may
be computed from any resolution of $A$ by objects $C^i\in\mathbf M$.}

Indeed, such a resolution exists by (i), and computes the $R^pF(A)$ because
the lemma makes every $C^i$ $F$-acyclic.

\TohokuContinuation{159}
